\documentclass[blue,iPad,cn]{elegantnote}

\title{\bfseries 中华乐教导引笔记}

\author{\href{www.liyunzhe.cn}{\itshape 李蕴哲} \\
		Wuhan University}	%\thanks{Elegant\LaTeX{} 其他模板下载地址：\href{https://ddswhu.me/resource/}{https://ddswhu.me/resource/}} }
\date{\today}

\usepackage{listings} 
\lstset{language=[LaTeX]{TeX},basicstyle=\footnotesize\ttfamily}


\begin{document}
{\color{ecolor}{\maketitle}}
% logo
%\centerline{\includegraphics[width=0.25\textwidth]{logo.png}}

\section{引言}

我们每个人是否是充实而快乐？

我们的出路和福祉何在？

山川异域，风月同天。而我们为什么只有“武汉加油”。

\subsection{日本送中国物资寄语}

\textbf{山川异域，风月同天。}这句话具有佛教色彩，但是也同样有东方美学，东方式的情感，体现了东方的诗意思维。\textbf{岂曰无衣，与子同裳}，来自诗经。这曾经是中国人很熟悉的表达方法，但是近几十年在国人的表达中渐渐疏离了。表达可以有多种方式，每一种表达都有使用的地方。但是这种古典失落的现象要重视。能雅能俗，能输能白。

\subsection{出路和福祉}

工具理性，科技思维是一种很好的思维。但是人的思维是要理性和感性并重的。

人的思维四个维度：

\begin{itemize}
	\item 智商。现在人基本不缺，甚至会自以为是。
	\item 情商。现在很多人已经在关注情商了，情商的提高做到三个方面：
	\begin{itemize}
		\item 努力突破自我中心，不要过于自我为是
		\item 多做换位思考
		\item 追求悉心体贴
	\end{itemize}
	\item 灵韵。当我们对美和艺术钟爱、在行的时候，灵韵就到了一个较好的状态。
	\item 情怀。东方文化首重情怀，情怀无价。
\end{itemize}

\subsection{真的快乐吗？}

中国文化倡导众美合一。纷吾既有此内美兮。但是也要发现自己的缺点。

重视自己内心的灵性，个体是谦和的，是卑微的，但是自己在气象上可以做到卓尔不群。要更加具有文化人的自觉，规避现代人的傲慢。实现宜古宜今的趣味与风雅。

\subsection{传统文化与文化传统}

文化传统就是一种文化的精气神。传统文化要有扬弃的态度。中华文化传统有雅俗的方面。文化传统的价值怎么估量都不过分，需要一个族群有价值的自觉，要敬畏文化传统。中华的文化传统在近一百多年有不同的际遇，这是一个沉重的话题。

\textbf{党和国家：不忘本来，吸收外来，面向未来。}如何践行，需要我们的考量。文化传统涉及到文化立场。每个人都有文化立场。文化立场，文化格调，文化风骨，是每个文化人都要思考的话题。

很多学者有过思考：

\begin{itemize}
	\item 饶宗颐：万古不磨意，中流自在心。
	\item 叶嘉莹：书生报国成何计，难忘诗骚李杜魂。
	\item 黄会林：大道本无我，青春长与君。
\end{itemize}

通过对文化个案的悉心领会去感受中华文化的本体。散为群星，聚为满月。中华文化就是月映万川，是具体性和独特性的统一。

\subsection{诗乐源流}

在中华文化的观念中，诗乐同源，古乐璀璨。古乐是上古时期的文化样态，常用来祭祀。诗经可以入乐歌唱，如今只有文本，甚至有些只有笙诗。只有到古代的情景下才能将内涵发扬出来。

《关雎》经学家的解读为赞颂后妃之德。在古乐场景下解读，不仅仅是表达对女子的爱慕。《关雎》其实是闹洞房，众人齐唱的歌。既有爱的热烈情思的流露，又有一种揶揄的色彩。

《蒹葭》、《野有蔓草》\footnote{体现了遇到什么美好就是什么的美好祝愿}都是先民的时代，传达的一种美好的祝愿。诗经的篇章中体现出的生命情调，灵感悟性，精神的丰腴。要融通传统与现代，贯通古今融汇有价值的美好的元素。王冲说过“善才有深浅，无有古今”。“文有真伪，武有古心”，要领会古今的文化。中华文化一通百通，是一种辩证统一。

到了汉代，有所分化，但是还有很多融合。诗乐交融是中华文化的一个很重要的色彩。词的本质也是音乐文学。晚唐词都是月前花下，入乐歌唱的。研读文本的时候，要注重文艺语境和文艺特征。

\subsection{中国文化中诗的内涵}

诗者，寺也。诗是把人内心的情愫以语言的形式展现出来。诗，源于情感的悦动。叶嘉莹认为，诗是兴发感动的载体。

诗乃天地之心。每个人喜欢的诗是不一样的。除了很纯美的诗，还有很朴素的诗。这都无妨诗的样貌。因为诗是天地之心，天地是多元的。

不同样态的中国诗：

\begin{itemize}
	\item 中国古代的抒情诗词，比如李商隐的作品《春雨》：	“怅卧新春白袷衣，白门寥落意多违。红楼隔雨相望冷，珠箔飘灯独自归。远路应悲春晼晚，残霄犹得梦依稀。玉珰缄札何由达，万里云罗一雁飞。”文饰之悲，沧海月明珠有泪。李商隐的诗，体现了抒情诗的极致。李清照的词，体现了贵族女性的才华，思维的灵妙。《一剪梅》婉约令词的巅峰状态。中国女性的柔善，两方着笔，“此情无计可消除，才下眉头，却上心头”。王菲演绎成了流行歌。纳兰容若又翻出新意。
	\item 很沉重的诗，比如李绅的《悯农》。杜甫的《又呈吴郎》：“堂前扑枣任西邻，无食无儿一妇人。不为困穷宁有此？只缘恐惧转须亲。即防远客虽多事，便插疏篱却甚真。已诉征求贫到骨，正思戎马泪盈巾。”杜甫的柔善，深沉的人文情怀，能忍常人之不能忍，更加的体恤众生\footnote{李白，不能忍常人之能忍，因此更加天马行空}。清人卢德水说：“杜诗温柔敦厚，其慈祥恺悌之衷，往往溢于言表。如此章，极煦育邻妇，又出脱邻妇；欲开导吴郎，又回护吴郎。八句中，百种千层，莫非仁音，所谓仁义之人其音蔼如也。”
\end{itemize}

数点梅花天地心。意象不再是自然之物了，具有更独特的传情达意的功能。

\begin{example}
	红豆意象。盛唐第一诗人王维的《相思》，红豆便成为了相思子，既浓且艳的美好内涵，再加上悲戚的意味。个体的情感依托更多的社会环境，因此后来更多的表现家国情怀。作者之意未必然，读者之意何必不然。中国文人没有不写红豆的。红楼梦里面宝玉的《红豆曲》。
\end{example}

\begin{example}
	梅花意象。梅花塑造了中国人的文化内涵，是中国的国花之一。孙中山先生说：“唯有梅花真国魂”。陆游的《卜算子咏梅》体现了中国文人的风骨，“零落成泥碾作尘，只有香如故”。对梅花的精神的赞誉。王冕的《墨梅诗》：“不要人夸好颜色，只留清气满乾坤。”高崎：“雪满山中高士卧，月明林下美人来。”“疏影横斜水清浅，暗香浮动月黄昏。”“十年不到香雪海，梅花亦我我亦梅。”
\end{example}

中国人人与自然交融，中华文化的华彩极为充浓。

“东风不来，三月的柳絮不飞。”柳絮有丰富的生命意义，又有漂泊、悲惨融入其中。中国人对诗的理解，呈现的诗意的思维和表达的精妙。

中国诗乐文化的璀璨成果：

\begin{itemize}
	\item 诗词。国诗的概念。
	\item 昆曲。
	\item 古琴。唯一可以称琴道的乐器。琴乃文器之心。
	\item 红楼梦。
\end{itemize}

\subsection{诗化的感性能力和理性能力}

雅人高致。古人强调，修得一颗明心。苏雪林，“人雅不关居有著，鸟鸣如唤克鹈鹕。”深情妙赏之人，往往有些许自恋，往往异常可爱。由于现代的逻辑思维，我们要规避逻辑的束缚，避免欲念的迷狂。中华文化的辩证，通达与和谐。

古人有一种更诗意的感觉。现在，雨水刚过，惊蛰将来，春气发扬。梅花已谢杏花新。古人认为这个是二十四方花性\footnote{每五天有一种新的花开放}的节点。

中国的诗乐塑造了柔善的人格样态，崇刚恋阴，婉约风致。

\begin{example}
	沈祖棻的代表作1932年的《浣溪沙》，张充和1943年的作品《临江仙》，体现了婉约风致。
\end{example}

期待这个课程能成为一段隽永的记忆！

\end{document}
